\documentclass[12pt]{article}
\usepackage{seantex}

% --- Document Info ---
\title{ÜS: Grundstrukturen, Woche 10}
\author{Sean Findlay}
\date{\today}

% --- Start Document ---
\begin{document}

\maketitle

Wann funktioniert die Methode, dass wir die Vereinigung als obere Schranke nehmen? Wenn $\mathcal{F}$ endlichen Charakter hat, dann funktioniert sie.

\section{Kardinalzahlarithmetik}

\begin{definition}{}{}
    Seien $\kappa, \lambda$ zwei Kardinalzahlen. Dann gilt:
    \begin{itemize}
        \item $\kappa + \lambda = |\{\kappa \times \{0\} \cup \lambda \times \{1\}\}|$
        \item $\kappa \cdot \lambda = |\kappa \times \lambda|$
        \item $\kappa^\lambda = |{}^\lambda \kappa|$
    \end{itemize}
\end{definition}

\begin{exercise}{}{}
    Zeige mit Kardinalzahlarithmetik:
    \begin{enumerate}
        \item $1 + 2 = 3$
        \item $1 \cdot 2 = 2$
        \item $2^2 = 4$
    \end{enumerate}
\end{exercise}

\begin{proofbox}
    \begin{enumerate}
        \item
        \[
        \begin{aligned}
            1 + 2 &= |\{1\times \{0\} \cup 2 \times \{1\}\}| = |\{\{\emptyset\}\times \{\emptyset\} \cup \{\emptyset, \{\emptyset\}\} \times \{\{\emptyset\}\}\}| \\
            &= |\{(\{\emptyset\}, 0), (\emptyset, 1), (\{\emptyset\}, 1)\}| = 3
        \end{aligned}
        \]

        \item
        \[
        \begin{aligned}
            1 \cdot 2 &= |1 \times 2| = |\{\emptyset\} \times \{\emptyset, \{\emptyset\}\}| = |\{ (\emptyset, \emptyset), (\emptyset, \{\emptyset\}) \}| = 2
        \end{aligned}
        \]

        \item
        \[
        \begin{aligned}
            2^2 &= |{}^2 2| = |\{f: 2 \rightarrow 2, f \text{ Funktion}\}| = |\{f: \{\emptyset, \{\emptyset\}\} \rightarrow \{\emptyset, \{\emptyset\}\}, f \text{ Funktion}\}| \\
            &= |\{ f_1 = \emptyset, f_2 = \{\emptyset\}, f_3 = \mathrm{id}, (f_4(\emptyset) = \{\emptyset\}, f_4(\{\emptyset\}) = \emptyset) \}| = 4
        \end{aligned}
        \]
    \end{enumerate}
\end{proofbox}

\begin{exercise}{}{}
    Seien $X$ eine endliche Menge. Zeige: $|X| = |X \setminus \{x\}| + 1 \ \forall x \in X$ falls $X \neq \emptyset$.
\end{exercise}

\begin{proofbox}
    Es gibt $f: n \rightarrow X$ Bijektion mit $n \geq 0,\ n \in \omega$ da $x \neq \emptyset$. Es existiert $n'$, sodass $n = s(n') = n' \cup \{n'\}$. Das heisst, es existiert $m$ sodass $f(m) = x$. Dann existiert eine Bijektion $\sigma: n \rightarrow n$ sodass $m$ und $n$ vertauscht werden.
    
    Sei dann $g = f \circ \sigma: n \rightarrow x \iff n' \cup \{n'\} \rightarrow (X \setminus \{x\}) \cup \{x\}$. Das heisst es gibt eine Bijektion $n' \rightarrow X \setminus \{x\}$. Es gilt $|x| = n = s(n') = n' + 1 = |X \setminus \{x\}| + 1$.
\end{proofbox}

\begin{exercise}{}{}
    Seien $X, Y$ endliche Mengen. Es gilt: $|X \cup Y| + |X \cap Y| = |X |+ |Y|$. Zeige das folgendes gilt:
    $$
        |X \cup Y| \leq |X| + |Y|
    $$

    mit Gleichheit genau dann, wenn $X \cap Y = \emptyset$.
\end{exercise}

\begin{proofbox}
    $X \cap Y \neq \emptyset \implies |X \cap Y| = k > 0 \implies k = s(k')$ für ein $k' \in \omega$.
    
    $\implies |X \cup Y| + k = |X| + |Y|$

    $\implies |X \cup Y| + k' + 1 = |X| + |Y|$

    $\implies |X \cup Y| + k' < |X \cup Y| + k' + 1 = |X| + |Y|$

    Wenn $X \cap Y = \emptyset$, dann ist $k = |\emptyset| = 0$. Also haben wir Gleichheit.
\end{proofbox}

\begin{exercise}{}{}
    Sei $X$ eine endliche Menge, sei $n = |X|$, $\mathrm{Perm}(X) = \{ f \in {}^X X\ |\ f \text{ ist bijektiv} \}$. Was ist $|\mathrm{Perm}(X)|$?
\end{exercise}

\begin{proofbox}
    Beweis per Induktion.

    Induktionshypothese: $|\mathrm{Perm}(X)| = n!$.

    \begin{itemize}
        \item \underline{Induktionsverankerung}: $n = 0$, dann ist $\mathrm{Perm}(X) = \{\emptyset\}$. Wir haben $|\mathrm{Perm}(X)| = 1 = 0!$.
        \item \underline{Induktionsschritt}: Es gilt $n \geq 1$, also existiert $n'$, sodass $n = s(n')$. Wir wissen aus der Induktionsannahme, dass $\mathrm{Perm}(X_{n'}) = n'!$ gilt. Sei $f: n' \cup \{n'\} \rightarrow n' \cup \{n'\}$ eine Bijektion.
        \begin{itemize}
            \item \underline{Fall 1: $f(n') = n'$}. Dann ist $f|_{n'}: n' \rightarrow n'$ eine Bijektion.
            \item \underline{Fall 2: $f(n') = m \in n'$}. Dann gibt es $\sigma: n \rightarrow n$ sodass $m$ und $n'$ vertauscht werden. Dann gilt $(f \circ \sigma)|_{n'} : n' \rightarrow n'$ eine Bijektion.
        \end{itemize}
    \end{itemize}
\end{proofbox}

\end{document}